\documentclass[10pt,a4paper]{article}
\usepackage{xeCJK}
\setCJKmainfont{Yu Mincho}
\setCJKsansfont{Yu Gothic}
\setCJKmonofont{Yu Gothic}
\usepackage{amsmath,amssymb,bm}
\usepackage{enumitem}
\usepackage[margin=25mm]{geometry}
\usepackage{hyperref}

\title{d'Haultfoeuille (2010) 注釈付き日本語読解ノート}
\author{読解対象：Xavier d'Haultfoeuille, Journal of Econometrics 154 (2010), 1--15}
\date{}

\newcommand{\indep}{\mathrel{\perp\!\!\!\perp}}
\newcommand{\E}{\mathbb{E}}
\newcommand{\Prb}{\mathbb{P}}
\newcommand{\note}[1]{\par\smallskip\noindent\textbf{注釈：}#1\par\smallskip}
\newcommand{\connect}[1]{\par\smallskip\noindent\textbf{欠測IVメモへの接続：}#1\par\smallskip}
\newcommand{\cue}[1]{\par\smallskip\noindent\textbf{原文の短い手がかり：}#1\par\smallskip}

\begin{document}
\maketitle

\section*{書誌情報と読解方針}

本ノートは，Xavier d'Haultfoeuille による ``A new instrumental method for dealing with endogenous selection'' を，研究メモとして日本語で順に読み下すための注釈付き資料である。対象論文は \textit{Journal of Econometrics}, 154(1), 1--15 に掲載された。

著作権上，本ノートは原論文の逐語的な全文翻訳ではなく，論文全体の議論を順番どおり追えるようにした詳細な日本語パラフレーズである。原文確認のための英語は短い専門句に限る。数式・仮定・定理は，読解上必要な範囲で再構成して説明する。

\subsection*{用語対応表}

\begin{center}
\begin{tabular}{ll}
endogenous selection & 内生的選択，選択バイアス \\
selection dummy & 選択ダミー，観測ダミー \\
outcome / dependent variable & アウトカム，被説明変数 \\
instrument & 操作変数，欠測IV，shadow variable \\
nonignorable nonresponse & 無視できない無回答，非無作為欠測 \\
completeness & 完備性 \\
inverse problem & 逆問題 \\
sharp bounds & 鋭い境界 \\
inverse probability weighting & 逆確率重み付け \\
\end{tabular}
\end{center}

\connect{本プロジェクトの \texttt{manuscript/main.tex} では，d'Haultfoeuille 型の欠測IVをベクトル値欠測変数と潜在因子モデルへ拡張する。本ノートでは，その基礎となるスカラー版の識別論，完備性，境界識別，推定法を整理する。}

\section{Abstract}

\cue{endogenous selection / conditional on the outcome / completeness}

要旨は，通常の操作変数法とは異なる発想を提示する。標準的なIV戦略では，操作変数がアウトカムから独立である，あるいはアウトカム方程式に直接入らないことを重視する。しかし選択がアウトカム自体に強く依存している場合，この発想では有効な変数を見つけにくい。そこで論文は，アウトカムを条件づけた後で，操作変数と選択変数が独立になる，という別の除外制約を使う。

この枠組みは，無視できない無回答，欠落共変量を含む二値モデル，観測されない部門を持つRoyモデルなどに向く。点識別は，操作変数とアウトカムの関係が十分に強いことを表す完備性条件によって得られる。完全な条件付き独立が疑わしい場合でも，選択確率がアウトカムや操作変数に単調に反応するという弱い仮定の下で，関心パラメータの鋭い境界を導く。最後に，推定法とフランス小学校の留年効果の応用を扱う。

\note{要旨の中心は，欠測確率 $\Prb(D=1\mid Y)$ を回復すれば母集団分布を逆確率重み付けで戻せる，という点である。通常のMARでは $\Prb(D=1\mid X,Z)$ を観測データから直接推定するが，本論文では選択が $Y$ に依存するため，その確率を積分方程式から逆算する。}

\section{1. Introduction}

導入部は，欠測値・選択・反実仮想の非観測がマイクロデータで一般的であり，観測集団だけで推論すると通常は不一致になる，という問題設定から始まる。追加仮定がない場合，Manski 型の部分識別のようにパラメータの範囲しか得られない。

点識別の既存アプローチとして，第一にMARまたはunconfoundednessがある。これは観測共変量を条件づけると選択とアウトカムが独立になるという仮定であり，選択とアウトカムの未観測の相関を排除するため，しばしば強すぎる。

第二に，選択には効くがアウトカムには直接効かない通常型の操作変数がある。標本選択モデルや処置効果モデルで用いられるが，これだけではアウトカム分布全体を一般に点識別しない。また，選択がアウトカムに強く依存する場面では，選択だけに効く変数を見つけることが難しい。

第三に，線形性や特別な回帰変数の存在など，関数形制約を使う方法がある。Chamberlain の identification at infinity や Lewbel の special regressor アプローチが例である。ただし，広いサポートや選択確率の端での挙動など，別の強い条件が必要になる。

本論文はこれらと異なり，選択ダミーと操作変数がアウトカムを条件づけると独立になるという仮定を使う。これは無回答の文脈では先行研究があるが，本論文の焦点はノンパラメトリック識別である。切断や無回答がアウトカムに直接依存する場合，選択だけに効く変数ではなく，アウトカムに効く変数をIVとして使える点が重要である。

通常のIV回帰と同じく，単なる除外制約だけでは不十分であり，IVとアウトカムの関係が十分に豊かでなければならない。本論文ではこのランク条件を完備性で表す。完備性と条件付き独立により，データの結合分布をノンパラメトリックに識別できる。

識別の鍵は，アウトカム条件付きの選択確率 $P(y)=\Prb(D=1\mid Y=y)$ を回復することにある。MARで傾向スコアを推定するのに似ているが，MARでは傾向スコアが観測共変量から直接推定できるのに対し，本論文では $Y$ が欠測するため，$P$ は逆問題として求める。

論文はさらに，操作変数の全分布ではなく一部のモーメントしか使えない場合にも，選択確率をパラメトリックに置けば識別できることを示す。これは補助情報やsurvey calibrationの発想に近く，Nevo 型の結果を内生的選択へ拡張する。

除外制約は応用上強く見えるが，MARと違って検定可能な含意を持つ。また，除外制約が完全には成り立たない場合でも，アウトカムと操作変数が選択に単調に影響するなら，関心パラメータの有限で鋭い境界が得られる。従属変数が非有界でも，平均などのパラメータにコンパクトな区間を与える点が特徴である。

推定面では，離散アウトカムやパラメトリックモデルではGMMを用いる。連続アウトカムのノンパラメトリック設定では，未知関数が線形逆問題の解になるため，Tikhonov正則化を使う。識別・推定に続いて，Monte Carloで有限標本性能を調べ，フランスの小学5年生の留年効果へ応用する。

\connect{導入部の貢献は，\texttt{main.tex} の「欠測IV/shadow variableは，選択に効く通常IVではなく，欠測変数に情報を持つ常時観測変数でよい」という主張の直接的な土台である。}

\section{2. Identification}

\subsection{2.1 設定と主結果}

記号は，$D$ が選択または観測ダミー，$Y$ がアウトカム，$Z$ が操作変数である。最初の二つの仮定はデータの観測構造を定める。

\begin{description}[leftmargin=3.2em]
\item[仮定1] $D$ は常に観測され，$D=1$ のとき $(Y,Z)$ が観測される。$D=0$ のとき $Y$ は観測されない。
\item[仮定2] $Z$ の分布は識別されている。
\end{description}

仮定1と2は，$Y$ だけが欠測する項目無回答や選択問題では自然である。単位無回答では $(Y,Z)$ がともに欠測しうるため，$Z$ の周辺分布をrefreshment sample，センサス，行政データなどから得る必要がある。後の2.4節では，$Z$ の全分布ではなく一部のモーメントだけで済ませるパラメトリック版を扱う。

仮定1と2だけでは $(D,Y,Z)$ の結合分布は点識別されない。選択が $Y$ に依存するなら，外生的選択は破綻する。そこで本論文は，$Y$ と関係するが，$Y$ を条件づけると $D$ に直接関係しない変数 $Z$ を考える。

\begin{description}[leftmargin=3.2em]
\item[仮定3] $D\indep Z\mid Y$。
\end{description}

これは，選択方程式が $Y$ に依存し，$Y$ は $D=0$ で欠測するため直接推定できないが，$Y$ に影響する $Z$ を使えば通常のIV回帰に似た形で選択方程式を識別できる，という発想である。構造的に
\[
Y=\varphi(Z,\varepsilon),\qquad D=\psi(Y,\eta)
\]
と書け，$\eta\indep (Z,\varepsilon)$ なら仮定3が成立する。これが命題2.1である。

\begin{description}[leftmargin=3.2em]
\item[命題2.1] 上の非パラメトリックシステムで $\eta$ が $(Z,\varepsilon)$ から独立なら，$D\indep Z\mid Y$ が成立する。
\end{description}

\note{この命題は，欠測IVのDAG的直観を与える。$Z\to Y\to D$ という経路はあるが，$Y$ を条件づけると $Z$ から $D$ への追加経路が消える。}

次に必要なのが完備性である。$B$ を，$g(Y)$ が下に有界で可積分である関数の集合とする。

\begin{description}[leftmargin=3.2em]
\item[仮定4] 任意の $g\in B$ について，
\[
\E\{g(Y)\mid Z\}=0\ \text{a.s.}\quad\Rightarrow\quad g(Y)=0\ \text{a.s.}
\]
が成り立つ。
\end{description}

これは通常の完備性よりは弱いが，有界完備性よりは強い。離散サポートの場合，$Y$ と $Z$ の条件付き分布行列がフルランクであることに対応する。したがって $Z$ のサポートは少なくとも $Y$ と同じくらい豊かで，条件付き分布が十分に線形独立でなければならない。

\begin{description}[leftmargin=3.2em]
\item[命題2.2] $Y=\mu\{\nu(Z)+\varepsilon\}$，$Z\indep\varepsilon$，$\nu(Z)$ が実数全体に十分なサポートを持ち，$\varepsilon$ の密度と特性関数に正則性があれば，$Y$ は $Z$ に対して $B$-completeである。
\end{description}

命題2.2は，順序選択モデル，Tobit，duration modelなどを含む広い非線形変換モデルで完備性が成り立つ十分条件を示す。重要なのは，完備性が単なる相関の強さではなく，条件付き期待作用素の核が自明であることを要求する点である。

\begin{description}[leftmargin=3.2em]
\item[仮定5] $P(Y)=\Prb(D=1\mid Y)>0$ がほとんど確実に成り立つ。
\end{description}

これは処置効果文献の共通サポート条件に相当する。$D$ が $Y$ の決定論的な切断関数である場合には破れる。ランダム切断でも，切断閾値が $Y$ の下限を常に上回るなら破れる。

\begin{description}[leftmargin=3.2em]
\item[定理2.3] 仮定1--5の下で，$(D,Y,Z)$ の分布は識別される。
\end{description}

証明の中心は，未知関数 $Q$ に関する方程式
\[
\E\left\{\frac{D}{Q(Y)}\mid Z\right\}=1
\]
である。真の $P(Y)=\Prb(D=1\mid Y)$ はこの方程式を満たす。もし別の解 $Q$ があれば，$P/Q-1$ の条件付き期待が $Z$ の下でゼロになり，完備性によって $P=Q$ が従う。よって選択確率が識別される。選択確率が分かれば，
\[
f_{Y,Z}(y,z)=\frac{f_{D,Y,Z}(1,y,z)}{P(y)}
\]
の形で母集団分布を復元できる。

\connect{\texttt{main.tex} の主定理では，スカラー $Y$ をベクトル $\bm Y$ に置き換え，$R\indep V\mid \bm Y$ と vector completeness によって同じ方程式 $\E\{R/\pi(\bm Y)\mid V\}=1$ を使う。}

\subsubsection{例1：無視できない無回答}

第一の例では，$Y$ は欠測しうるアウトカム，$Z$ は全員について分布が分かる補助変数である。選択が $Y$ に依存してMARでない場合でも，$D\indep Z\mid Y$ と完備性があれば，$Y$ の分布を識別できる。単位無回答で $Z$ 自体が本標本では欠測しても，外部データで $Z$ の分布が分かれば仮定2を満たせる。

\subsubsection{例2：観測されない部門を持つRoyモデル}

第二の例は，労働市場の部門選択や賃金観測の問題である。ある部門の賃金 $Y$ が選択された人にしか観測されないとき，選択は潜在賃金に依存しうる。通常IVとしては選択だけに効く変数を探すが，本論文の方法では賃金に効く変数を使える。たとえば教育や経験のように賃金に関係する変数が，賃金を条件づけた後では選択に直接効かないなら，欠測IVとして使える。

\subsubsection{例3：一つの応答層だけからの標本}

第三の例は，二値応答モデルで一方の応答層だけが観測される場合である。たとえば $D=1$ の標本だけで $Y,Z$ が得られる状況でも，$Z$ の分布が別途識別され，$D\indep Z\mid Y$ が成り立てば，全体分布を復元できる。

\subsection{2.2 検定可能性}

この節は，仮定3が単なる信念ではなく，観測分布上の含意を持つことを示す。任意の候補関数 $Q$ について
\[
\E\left\{\frac{D}{Q(Y)}-1\mid Z\right\}=0
\]
が観測データから評価できる。もし $(0,1]$ に値を取る解 $Q$ が存在しなければ，仮定3は棄却される。

\begin{description}[leftmargin=3.2em]
\item[定理2.4] 仮定1,2,5の下で，方程式に $(0,1]$ 値の解が存在しない場合，かつその場合に限り，仮定3は棄却できる。
\end{description}

この結果は，条件付き独立を完全に証明できるという意味ではない。解が存在すれば，仮定3と整合的な結合分布を構成できる，という意味である。離散 $Y$ の場合はGMMで検定できる。連続 $Y$ の場合は，ノンパラメトリック逆問題の残差に基づく検定が考えられる。

\note{MARは一般に観測データだけでは検定しにくい。これに対して本論文の欠測IV仮定は，$Z$ の分布情報と完全ケースの条件付き分布を使って反証可能な含意を持つ。}

\subsection{2.3 条件付き独立なしの集合識別}

この節では，$D\indep Z\mid Y$ が疑わしい場合に，より弱い単調性で何が言えるかを扱う。$Z$ だけでなく，分布が識別され $Z$ から独立な別の操作変数 $\widetilde Z$ も導入される。

\begin{description}[leftmargin=3.2em]
\item[仮定3'] $z\mapsto \Prb(D=1\mid Y,Z=z)$ はほとんど確実に単調増加である。
\item[仮定6] $y\mapsto \Prb(D=1\mid Y=y,\widetilde Z)$ はほとんど確実に単調増加である。
\end{description}

仮定3'は，$Y$ を固定したとき，$Z$ が高いほど選択されやすいという条件であり，仮定3の独立性を単調性へ弱めたものである。仮定6は，$\widetilde Z$ を固定したとき，$Y$ が高いほど選択されやすいという条件であり，MARを単調性へ弱める。

定理2.5は，$h(Y)$ の平均の上限・下限を与える。$h$ が増加関数である場合，単調性と共分散の符号を使って，不観測部分の平均を完全ケースや補助変数で挟む。

\begin{description}[leftmargin=3.2em]
\item[定理2.5] 仮定1と正の選択確率の下で，仮定6から上限が得られ，仮定3'と適切な関数クラス条件から下限が得られる。境界は鋭い。
\end{description}

境界が鋭いとは，その範囲内の値を生むようなデータ生成過程が仮定と観測分布に整合的に存在する，という意味である。したがって，これは保守的な便宜的区間ではなく，仮定から論理的に導かれる最良の区間である。

\begin{description}[leftmargin=3.2em]
\item[命題2.6] 関数 $h$ が定理2.5の必要な関数クラスに入るための実用的な条件を与える。条件付き分布の一次確率優越や線形仕様の検定が関係する。
\end{description}

\note{この節は，完全な欠測IV除外制約が成り立たない場合の逃げ道である。研究上は，主分析を点識別，頑健性分析を単調性境界で提示する構成が考えられる。}

\subsection{2.4 パラメトリック識別}

ノンパラメトリック識別は汎関数方程式の一意解に依存するため，実装では次元の呪いや不安定性が問題になる。また，仮定2のように $Z$ の全分布が分かるとは限らない。そこで，$Z$ の一部のモーメント，たとえば $\E(W)$ だけが既知で，選択確率にパラメトリック制約を置く場合を考える。

共変量 $X$ を明示し，$V=(X',Y')'$，$W=(X',Z')'$ とする。選択確率は
\[
\Prb(D=1\mid V)=F(V'\beta_0)
\]
と置かれる。$F$ は既知の単調なリンク関数で，ロジットやプロビットが典型例である。

\begin{description}[leftmargin=3.2em]
\item[仮定2'] $\E(W)$ が既知であり，選択確率は既知リンクを持つ単一指数モデルであり，$V'\delta=0$ が完全ケースで成立すれば $\delta=0$ である。
\item[仮定3'] $D\indep Z\mid V$。
\item[仮定4'] 局所識別のためのランク条件。
\item[仮定4''] $Z$ の完全ケース条件付き平均が $X,Y$ に線形で，$Y$ 側の係数行列がフルランクである。
\end{description}

\begin{description}[leftmargin=3.2em]
\item[定理2.7] 仮定1,2',3'の下で，$\beta_0$ は仮定4'が成り立つ場合に局所識別される。仮定4''が成り立つなら大域識別される。
\end{description}

この節の含意は，外部データから $Z$ の全分布が得られなくても，補助モーメントと選択モデルのパラメトリック形を組み合わせれば識別できるということである。これはsurvey calibrationや補助情報を用いた重み付けの発想と近い。

\section{3. Estimation}

推定節は，選択確率 $P$ またはその逆数 $f=1/P$ をどのように推定するかを扱う。標本は独立同分布で，$Y^*=DY$ として観測される。

\begin{description}[leftmargin=3.2em]
\item[仮定7] $((D_i,X_i,Y_i^*,Z_i))_{i=1}^n$ が独立コピーとして観測される。
\end{description}

\subsection{3.1 パラメトリック推定}

$Y$ が有限個の値を取る場合，定理2.3の方程式を有限次元のモーメント条件として書ける。したがって，各点の選択確率 $P(y_k)$ をGMMで推定できる。

2.4節のパラメトリック選択確率を用いる場合も，
\[
\E\left[\left\{\frac{D}{F(V'\beta_0)}-1\right\}W\right]=0
\]
というモーメント条件が得られる。仮定4''の下ではこのGMM目的関数が $\beta_0$ を大域的に識別する。

\subsection{3.2 ノンパラメトリック推定}

連続 $Y$ でパラメトリック制約を避けたい場合，未知関数 $f=1/P$ を条件付きモーメントから推定する。説明を簡単にするため，$X$ はなく，$(Y,Z)\in[0,1]^2$ とする。作用素
\[
T\phi(z)=\E\{D\phi(Y^*)\mid Z=z\}
\]
を定義すると，識別方程式は
\[
Tf=1
\]
である。

この逆問題はill-posedなので，経験作用素 $\widehat T$ をカーネルで推定し，Tikhonov正則化によって
\[
\widehat f\in\arg\min_{\phi\in D_M}\|\widehat T\phi-1\|^2+\alpha_n\|\phi\|^2
\]
を解く。$D_M$ は $1\leq \phi\leq M$ を満たす関数集合であり，$f=1/P$ が1以上であることと対応している。

\begin{description}[leftmargin=3.2em]
\item[仮定8] $f$ は $D_M$ に属し，$(Y,Z)$ の密度と周辺密度が適切に有界である。
\item[仮定9] カーネルは非負で積分1，一次モーメント0を満たす。
\item[仮定10] バンド幅と正則化パラメータが適切な速度で0に近づく。
\item[定理3.1] 仮定3,4,7--10の下で，$\widehat f$ は $L^2$ 平均で $f$ に一致する。
\end{description}

次に，任意の汎関数 $\theta=\E\{g(Y,Z)\}$ を逆確率重み付き推定量
\[
\widehat\theta=\frac1n\sum_{i=1}^n D_i\widehat f_{-i}(Y_i^*)g(Y_i^*,Z_i)
\]
で推定する。leave-one-out の $\widehat f_{-i}$ を使うのは第一段階推定との依存を制御するためである。

\begin{description}[leftmargin=3.2em]
\item[系3.2] 定理3.1の条件下で，$\widehat\theta$ は $\theta$ に一致する。
\end{description}

\note{ここでの推定は，MARのIPWと形式は似ているが，重みを直接推定するのではなく，欠測IV方程式を正則化して求める点が本質的に異なる。}

\section{4. Monte Carlo simulations}

Monte Carloでは，パラメトリック推定量とノンパラメトリック推定量の有限標本性能を比較する。データ生成過程は
\[
Y=\Lambda\{\Lambda^{-1}(Z)+\varepsilon\},\qquad
D=1\{P(Y)\geq \eta\}
\]
であり，$Z$ と $\eta$ は一様分布，$\varepsilon$ は正規分布で独立である。選択確率 $P(y)$ は応用例で推定される形に近くなるよう選ばれている。これにより，命題2.1と2.2から仮定3と4が満たされる。

比較される第一段階推定量は二つのノンパラメトリック推定量 $\widehat f_1,\widehat f_2$ と，柔軟なパラメトリック族に属する $\widehat f_3$ である。性能指標は $f$ のMISEと，$\theta=\E(Y)=1/2$ のIPW推定量のbiasとRMSEである。さらに，真の $f$ を使う実行不能推定量も比較される。

表1の結果は，サンプルが小さいとパラメトリック推定量は不安定になりうるが，サンプルが大きくなると，仕様誤差のバイアスより精度改善が勝ち，パラメトリック推定量が最も良くなることを示す。興味深い点として，第一段階でより良い $f$ 推定量が必ずしも第二段階の $\theta$ 推定量を良くするとは限らない。また，真の重みを使う推定量より推定重みを使う推定量の方が分散面で有利になりうる。

\note{これは傾向スコア推定で，真の傾向スコアを使うIPWより推定傾向スコアを使うIPWが効率的になる現象と類似する。欠測IVでも第一段階推定が単なる誤差源ではなく，モーメント制約を通じて分散を下げる場合がある。}

\section{5. Application}

\subsection{5.1 導入}

応用は，フランスの小学5年生における留年が翌年のテスト成績へ与える短期効果である。フランスでは留年が比較的多く，2002年時点で小学校中に少なくとも一度留年した生徒が相当数存在したが，その効果の厳密な測定は少なかった。

データは，1997年に小学校1年生へ入学した9641人を追跡するフランス教育省のパネルである。3年生開始時の標準化テストを $Z$，6年生開始時の2002年テストを $Y$，2003年テストを $Y_1$ とする。分析標本は，5年生に在籍し，2002年または2003年に6年生へ到達した生徒から，テスト欠測を除いた5467人である。

表2は，5年生で留年した生徒，6年生で留年した生徒，どちらも留年しなかった生徒の平均テストスコアを示す。5年生留年者は2002年の6年生テスト $Y$ が観測されず，2003年の $Y_1$ が観測される。6年生留年者は6年生テストを2回受けるため，$Y$ と $Y_1$ の両方が観測される。この構造が反実仮想を補う鍵になる。

関心パラメータは，5年生留年者に対する処置効果
\[
\Delta_{TT}=\E\{Y_1(0)-Y_1(1)\mid D=0\}
\]
である。$Y_1(0)$ は留年した場合の翌年成績，$Y_1(1)$ は進級していた場合の翌年成績である。$D=0$ では $Y_1(0)$ は観測されるが，$Y_1(1)$ は反実仮想である。

フランスでは留年決定を駆動する明確な外生ルールがないため，通常の処置IV戦略は難しい。そこで，留年者が仮に進級していた場合の成長が，ゼロ以上で，かつ同じ初期スコアを持ち進級後に6年生で留年した生徒の成長以下である，という境界条件を置く。これにより，$\Delta_{TT}$ は $E(Y\mid D=0)$ と $E\{h(Y)\mid D=0\}$ を含む上下限で表される。ここで $h(Y)$ は，6年生留年者に観測される進歩関数である。

しかし5年生留年者では $Y$ が欠測しているため，$E(Y\mid D=0)$ と $E\{h(Y)\mid D=0\}$ はそのままでは識別できない。そこで3年生テスト $Z$ を欠測IVとして使う。以下で，条件付き独立による点識別戦略と，単調性による境界戦略を比較する。

\subsection{5.2 実証戦略}

\subsubsection{5.2.1 条件付き独立戦略}

第一戦略では，5年生での留年 $D$ が，6年生開始時能力 $Y$ を条件づけると，3年生スコア $Z$ から独立であると仮定する。直観的には，教師が留年決定で現在能力を見ており，過去の3年生スコアが現在能力を通じてしか効かないなら成り立つ。

完備性も仮定すると，定理2.3により $P(Y)=\Prb(D=1\mid Y)$ が識別され，$E(Y\mid D=0)$ と $E\{h(Y)\mid D=0\}$ を完全ケースから逆確率重み付けで復元できる。論文はこの復元式から，$\Delta_{TT}$ の下限と上限を式(5.4)，(5.5)として表す。

実装では，$h$ をカーネル推定し，$P$ をMonte Carlo節と同じ柔軟なパラメトリック形で推定する。図2は，$h$ と $P$ の推定値を示す。$P$ は完全ケースにおけるテストスコア $Y$ に対して非線形に変化する。

\subsubsection{5.2.2 単調性戦略}

第二戦略では，条件付き独立は強すぎると見て，$Y$ と $Z$ の双方が進級確率に正の影響を持つという単調性だけを使う。$Y$ は5年生終了時能力の不完全な測定であり，3年生スコア $Z$ も現在能力や教師判断に追加情報を持つ可能性があるため，この戦略はより保守的である。

論文は，6年生に進級した生徒のうち，6年生で留年したかどうかを使って，$Y$ と $Z$ が進級に正に効くことをロジットで確認する。表3では，2002年の6年生スコア $Y$ の効果が大きく，3年生スコア $Z$ の効果も正である。これは単調性仮定を支持する証拠として使われる。

さらに，$h$ が増加関数であり，定理2.5の関数クラスに入ることを確認する。図2から $h$ の増加性が見られ，仕様検定でも必要な線形関係を棄却しない。これにより，$E\{h(Y)\mid D=0\}$ と $E(Y\mid D=0)$ の境界を定理2.5で得る。

\subsection{5.3 結果}

表4が主要結果である。完全な条件付き独立を仮定した場合，$\Delta_{TT}$ の推定識別区間は正の値だけを含み，留年の短期効果は不利な場合でも正であることを示唆する。

単調性だけに依拠すると結果はより曖昧になる。もし留年が3年生スコアだけに依存していたと考えるなら留年は有害に見えるが，これは現実的でない。実際には，$Y$ の方が $D$ への影響が大きいと考えられるため，最悪ケースでも真の効果はゼロ付近に近い可能性が高い。

論文の結論は，留年の短期効果にはかなり不確実性があるが，正である可能性が高い，というものである。これはシカゴの3年生を対象にしたJacob and Lefgrenの結果と整合的である一方，6年生についての結果より楽観的である。国ごとの留年決定ルールの違い，教師・親の裁量，測定誤差，特定集団への体系的な有利不利が結果に影響しうる。

\note{応用の構造は，本プロジェクトの欠測IV設定と非常に近い。5年生留年者で $Y$ が欠測し，$Z$ は常時観測され，$D$ は $Y$ に依存する。ここで $Z$ を通常IVではなく，欠測アウトカム $Y$ に情報を持つshadow variableとして使う。}

\section{6. Conclusion}

結論部は，内生的選択への新しい操作変数アプローチをまとめる。識別の中心仮定は，通常の選択問題の仮定と対照的に，アウトカムを条件づけたときの操作変数と選択変数の独立性である。完備性の下で，一般的なノンパラメトリック点識別が得られる。

この枠組みは，観測されない部門を持つRoyモデル，無視できない無回答，一つの応答層だけが観測される二値モデルなど，幅広い選択問題に適用できる。条件付き独立が破れる場合でも，単調性の下で鋭く有限な境界を得られる。

今後の課題として二点が挙げられる。第一に，選択が実現アウトカムではなく予測アウトカムに依存する一般化Royモデルへ拡張できるかである。この場合，条件付き独立は崩れるが，モデル構造が点識別または集合識別の情報を与えるかもしれない。第二に，上限境界が抽象的なパラメータ集合上で得られているため，理論・実務の両面でより具体的な特徴づけが望まれる。

\section{Appendix proofs}

\subsection*{A.1 命題2.1の証明}

\paragraph{Original proof sequence (typeset).}
For fixed $(y,z)$, define the inverse-image sets
\[
 A_y=\{u:\psi(y,u)=1\},
 \qquad
 C_{y,z}=\{v:\varphi(z,v)=y\}.
\]
Then the proof proceeds as
\begin{align*}
 \Prb(D=1\mid Y=y,Z=z)
 &=\Prb(\eta\in A_y\mid Y=y,Z=z)\\
 &=\Prb(\eta\in A_y\mid \varepsilon\in C_{y,z},Z=z)\\
 &=\Prb(\eta\in A_y)\\
 &=\Prb(\eta\in A_y\mid Y=y)\\
 &=\Prb(D=1\mid Y=y).
\end{align*}
The middle equalities use $\eta\indep(Z,\varepsilon)$; hence
$D\indep Z\mid Y$.

\paragraph{日本語による論理の再確認.}
$D=\psi(Y,\eta)$ であるため，$Y=y$ を固定すると $D=1$ は
$\eta\in A_y$ と同じ事象である。また $Y=\varphi(Z,\varepsilon)$
なので，$(Y,Z)=(y,z)$ を固定することは $Z=z$ と
$\varepsilon\in C_{y,z}$ を固定することに対応する。
$\eta$ は $(Z,\varepsilon)$ から独立であるため，この条件づけは
$\eta$ の分布を変えない。最後に $Y$ も $(Z,\varepsilon)$ の関数で
あるから $\eta\indep Y$ であり，$Y$ のみを条件づけても同じ確率に
なる。したがって $Z$ は $Y$ を所与とした選択確率へ追加情報を
与えない。

\subsection*{A.2 命題2.2の証明}

証明は三段階である。第一に，誤差密度 $f_\varepsilon$ と安定分布密度を畳み込んだ関数が，上下から適切な多項式オーダーで有界であることを示す。第二に，$\E\{g(Y)\mid Z\}=0$ という条件を，$g\circ\mu$ とある密度の畳み込みがゼロであるという形へ変換する。第三に，Mattnerの位置族完備性定理を使い，特性関数が消えないことから $g\circ\mu=0$，したがって $g(Y)=0$ を得る。

\note{この証明は，完備性が「IVが強い」という直観を，畳み込み作用素の単射性として表現している。特性関数がゼロを持たないことは，deconvolutionで情報が完全には失われないことに対応する。}

\subsection*{A.3 定理2.3の証明}

\paragraph{Original proof sequence (typeset).}
Let $P(y)=\Prb(D=1\mid Y=y)$. Assumption 3 gives
\begin{align}
 \E\left\{\frac{D}{P(Y)}\middle|Z\right\}
 &=\E\left[\left.
   \frac{\E(D\mid Y,Z)}{P(Y)}\right|Z\right] \notag\\
 &=\E\left[\left.
   \frac{\E(D\mid Y)}{P(Y)}\right|Z\right]
 =1.
 \label{eq:annotated-a7}
\end{align}
For any proposed selection probability $Q:\mathcal Y\to(0,1]$,
the left-hand side of
\[
 \E\left\{\frac{D}{Q(Y)}-1\middle|Z\right\}=0
\]
is determined by the observed complete-case distribution and the
marginal distribution of $Z$. Thus every candidate for $P$ must solve
this equation.

Suppose $Q$ is a solution and put
\[
 g(Y)=\frac{P(Y)}{Q(Y)}-1.
\]
Because $P(Y)>0$ and $Q(Y)>0$, $g(Y)\geq-1$. Integrating the candidate
moment over $Z$ yields
\[
 \E\left\{\frac{P(Y)}{Q(Y)}\right\}
 =\E\left\{\frac{D}{Q(Y)}\right\}=1,
\]
and therefore
\[
 \E\{|g(Y)|\}
 \leq \E\left\{\frac{P(Y)}{Q(Y)}\right\}+1<\infty.
\]
Hence $g$ belongs to the admissible class $B$. Moreover,
\begin{align*}
 0
 &=\E\left\{\frac{D}{Q(Y)}-1\middle|Z\right\}\\
 &=\E\left\{\frac{P(Y)}{Q(Y)}-1\middle|Z\right\}
 =\E\{g(Y)\mid Z\}.
\end{align*}
Assumption 4 implies $g(Y)=0$ almost surely, so $Q(Y)=P(Y)$ almost
surely. The solution, and therefore the selection probability, is
unique.

Finally, the complete-case subdensity is observed:
\[
 f_{D,Y,Z}(1,y,z)=f_{Y,Z\mid D=1}(y,z)\Prb(D=1).
\]
Conditional independence gives
\[
 P(y)=\frac{f_{D,Y,Z}(1,y,z)}{f_{Y,Z}(y,z)},
\]
so that
\begin{align*}
 f_{Y,Z}(y,z)
 &=\frac{f_{D,Y,Z}(1,y,z)}{P(y)},\\
 f_{D,Y,Z}(0,y,z)
 &=\frac{1-P(y)}{P(y)}f_{D,Y,Z}(1,y,z).
\end{align*}
Thus the full distribution of $(D,Y,Z)$ is identified.

\paragraph{日本語による論理の再確認.}
観測可能性は，$D/Q(Y)$ が $D=1$ の標本でしか $Y$ を必要としない
ことから従う。一意性で直接完備性を適用する対象は $P-Q$ や
$1/P-1/Q$ ではなく $g=P/Q-1$ である。この選択により
$g\geq-1$ が自動的に得られ，候補モーメントから可積分性も確認
できる。したがって，原論文の通常の完備性より弱い
$B$-completenessで十分である。$P$ の識別後は，完全ケースの
subdensityを $P(y)$ で割って母集団の $(Y,Z)$ 分布を戻し，残りを
$(1-P)/P$ 倍して $D=0$ 側も復元する。

\subsection*{A.4 定理2.4の証明}

「もし」方向は明らかである。逆方向では，$(0,1]$ に値を取る解 $Q$ が存在するなら，その $Q$ を選択確率として持つ仮想的な結合密度を構成する。この結合密度は観測された完全ケース分布と $Z$ の周辺分布に整合し，かつ $D\indep Z\mid Y$ を満たす。したがって解が存在する限り，観測データだけでは仮定3を棄却できない。

\subsection*{A.5 定理2.5の証明}

証明の補題は，単調増加関数同士の共分散が非負であるという事実である。これを，選択確率と関心関数 $h(Y)$ の単調性に適用する。仮定6の下では完全ケース条件付き平均が母集団平均の上限を与え，仮定3'の下では逆重み付き量が下限を与える。鋭さは，定理2.4の構成と同様に，境界を達成する分布を作れることから示される。

\subsection*{A.6 命題2.6の証明}

この証明は，$h(Y)=\E\{\widetilde h(Z)\mid D=1,Y\}$ と書ける条件を，条件付き分布の一次確率優越に結びつける。$\widetilde h$ が増加関数なら，積分表現とFubiniの定理により，$Y$ が増えたとき条件付き分布が右へ移ることが $h$ の単調性を保証する。

\subsection*{A.7 定理2.7の証明}

パラメトリック選択確率では，真の $\beta_0$ が
\[
\E\left\{\frac{DW}{F(V'\beta_0)}\right\}=\E(W)
\]
を満たす。局所識別は，このモーメント写像の微分がフルランクであることと同値である。大域識別では，$Z$ の完全ケース条件付き平均が $X,Y$ に線形で，$Y$ 側の係数がフルランクであることを使い，任意の別解 $\beta$ が $V'(\beta_0-\beta)=0$ を意味し，仮定2'の識別条件から $\beta=\beta_0$ が従う。

\subsection*{A.8 定理3.1の証明}

証明は，Tikhonov正則化推定量の標準的な一貫性議論である。経験作用素 $\widehat T$ が真の作用素 $T$ に近づくこと，正則化項 $\alpha_n$ が十分ゆっくりゼロへ向かうこと，関数集合 $D_M$ が弱閉凸であることを使う。$T$ の連続性と弱逐次閉性により，任意の弱極限が方程式 $Tf=1$ を満たし，完備性から解が一意であるため $\widehat f\to f$ が得られる。

\subsection*{A.9 系3.2の証明}

IPW推定量の誤差を，第一段階の $\widehat f-f$ による項と，真の $f$ を使った標本平均の項に分ける。後者は大数の法則でゼロへ収束する。前者はCauchy--Schwarz不等式と定理3.1によってゼロへ収束する。したがって $\widehat\theta$ は一致する。

\section{Referencesの位置づけ}

参考文献は三つの流れに整理できる。第一に，Manski，Manski and Pepper，Little and Rubin など，欠測・選択・部分識別の基礎文献である。第二に，Newey and Powell，Blundell et al.，Hu and Schennach，Hall and Horowitz など，完備性・非パラメトリックIV・逆問題の文献である。第三に，Jacob and Lefgren，Troncin，Cosnefroy and Rocher など，留年効果の教育経済・教育学文献である。

\connect{本プロジェクトで次に読むべき文献は，欠測IV側では Wang, Shao and Kim (2014)，Zhao and Shao (2015)，Miao and Tchetgen Tchetgen 系列，潜在変数・測定誤差側では Hu and Schennach (2008) である。}

\section*{全体まとめ}

d'Haultfoeuille (2010) の方法は，非無作為欠測を「選択方程式の内生性」と見なし，アウトカムを条件づけた除外制約と完備性で選択確率を識別する。基本方程式は
\[
\E\left\{\frac{D}{P(Y)}\mid Z\right\}=1
\]
であり，これは欠測IV研究における中心的な積分方程式である。点識別が疑わしいときは単調性による鋭い境界へ移行でき，推定ではGMMまたは正則化逆問題を使う。応用では，5年生留年者に欠測する6年生能力を，3年生スコアをshadow variableとして補完し，留年効果の境界を求めている。

\end{document}
